\documentclass[11pt]{article}

\usepackage[margin=0.6in]{geometry}
\usepackage[T1]{fontenc}
\usepackage[utf8]{inputenc}
\usepackage{lmodern}
\usepackage{array}
\usepackage{booktabs}
\usepackage{hyperref}
\usepackage{parskip}
\usepackage{ragged2e}
\usepackage{setspace}
\usepackage[table]{xcolor}
\usepackage[natbibapa]{apacite}

\definecolor{PrimaryRed}{HTML}{8B1E1E}
\definecolor{SoftBlue}{HTML}{DDEEFF}
\definecolor{SoftGray}{HTML}{F4F4F4}

\hypersetup{colorlinks=true,linkcolor=black,urlcolor=black,citecolor=black}
\setlength{\parindent}{0pt}
\setstretch{1.08}
\bibliographystyle{apacite}
\newcommand{\reviewsection}[1]{\vspace{0.07in}\noindent{\normalsize\bfseries\color{PrimaryRed}#1}\par\vspace{0.02in}}

\begin{document}

\hypersetup{pageanchor=false}
\begin{titlepage}
\newgeometry{margin=0.7in}
\thispagestyle{empty}

\begin{center}
{\large University of Rochester\\}
{\normalsize Warner School of Education\\[0.15in]}

{\color{PrimaryRed}\rule{\textwidth}{1.4pt}}\\[0.12in]
{\Huge\bfseries Sensory Walk Experience\\[0.08in]
\huge Sensory Access, Communication,\\[0.08in]
and Agency in an Elementary\\[0.08in]
Computer Science Classroom}\\[0.18in]
{\Large EDE 448 --- Module 1}\\[0.12in]

\vfill

{\color{PrimaryRed}\rule{0.78\textwidth}{0.7pt}}\\[0.22in]

\colorbox{SoftBlue}{%
\parbox{0.88\textwidth}{%
\centering\small\vspace{0.02in}
\textbf{From Visible Behavior to Environmental Evidence}\\[0.05in]
A sensory-detective analysis of how visual structure, sound, movement, materials, transitions, communication options, and recovery supports shape children's access to participation.
}%
}
\end{center}

\vfill

\begin{center}
\renewcommand{\arraystretch}{1.18}
\setlength{\tabcolsep}{4pt}
\begin{tabular}{>{\bfseries\RaggedLeft}p{1.55in} p{4.95in}}
Student & Piter Zacari Garcia Bautista \\
Assignment & Sensory Walk Experience \\
Course & EDE 448 --- Communication and Positive Behavioral Supports for Students with Autism and Other Complex Needs \\
Setting & Pine Brook Elementary School IGNITE computer science classroom \\
Observation Context & Spring 2026 student-teaching placement \\
Interpretive Lenses & Sensory Detective; Michigan Sensory-Friendly Environment Checklist; Puzzle Plan and EQUITAS \\
Submission Date & July 2026 \\
\end{tabular}
\end{center}

\vfill

\begin{center}
\begin{minipage}{0.92\textwidth}
\centering\itshape
A sensory-friendly environment does not ask children to become less complex in order to belong. It changes the conditions so that more children can participate, communicate, regulate, and learn with dignity.\\
{\color{PrimaryRed}\rule{4in}{0.8pt}}
\end{minipage}
\end{center}

\vfill

\begin{center}
\begin{minipage}{0.95\textwidth}
\small
This paper examines sensory access as part of a broader communication ecosystem. The \textbf{Puzzle Plan} treats complex educational problems as interconnected systems that require multiple forms of evidence, community-centered interpretation, and iterative redesign. The \textbf{EQUITAS} framework asks whether learning environments provide equitable, culturally responsive, and accountable pathways for students to access instruction, communicate needs, demonstrate competence, and exercise agency. Together, these lenses shift attention away from deficit explanations of behavior and toward the relationship among the learner, the task, the environment, and the supports available.
\end{minipage}
\end{center}

\restoregeometry
\end{titlepage}
\hypersetup{pageanchor=true}
\pagenumbering{arabic}

\reviewsection{Purpose and Method}

For this sensory walk, I revisited the IGNITE computer science classroom at Pine Brook Elementary School through the perspective of a child entering the room, locating materials, listening to directions, participating in a hands-on lesson, communicating a need, taking a break, and transitioning out. I used the Guggenheim's \textit{Creating a Sensory-Friendly Environment: Becoming a Sensory Detective} resource and Michigan's \textit{Sensory-Friendly Environment Checklist}. The Sensory Detective resource emphasizes examining a visitor's experience from beginning to end rather than judging a setting by one isolated feature \citep{guggenheim_sensory}. The Michigan checklist adds practical questions about noise, lighting, temperature, crowds, alternative routes, sensory maps, and access to a quiet space \citep{mddi_checklist}.

I paid attention to sight, sound, touch, movement, body position, predictability, communication demands, and opportunities for regulation. My central question was whether the environment gave children enough information, choice, communication access, and recovery support to participate with dignity.

\reviewsection{Summary of the Experience}

The IGNITE classroom supports elementary students through computer science, digital literacy, robotics, engineering, and hands-on design. Students may use Chromebooks, block-coding programs, robots, building materials, craft supplies, and visual lesson resources. These activities are engaging and provide multiple pathways into learning, but they can also create several forms of sensory input at the same time. Screens add visual stimulation; devices and peers produce sound; students and materials move around the room; and building activities involve varied textures and fine-motor demands.

The classroom already included several practices that supported sensory access. During my placement, I observed picture-based choices, visual representations of expectations, step-by-step directions, clear start-and-stop signals, structured turn-taking, movement opportunities, and flexible ways for students to respond. These supports increased predictability and reduced the amount of spoken information a child had to hold in working memory.

One experience especially shaped my interpretation of sensory-friendliness. A student with complex support needs wanted to keep a laptop nearby while participating in a building activity. Instead of treating the laptop as a distraction or requiring the student to participate in only one approved way, I allowed the student to keep it nearby for comfort while building with blocks. The student remained connected to the activity. This showed me that sensory support is not always a separate room or a specialized object. Sometimes it is an adult recognizing a child's regulation strategy and preserving access without lowering the learning expectation.

\newpage
\reviewsection{Sensory Evidence}

\renewcommand{\arraystretch}{1.22}
\setlength{\tabcolsep}{5pt}
\rowcolors{2}{SoftGray}{white}
\begin{tabular}{>{\RaggedRight\arraybackslash}p{1.05in} >{\RaggedRight\arraybackslash}p{2.45in} >{\RaggedRight\arraybackslash}p{3.15in}}
\toprule
\rowcolor{SoftBlue}
\textbf{Area} & \textbf{Evidence from the classroom} & \textbf{Meaning for children with sensory needs} \\
\midrule
Visual structure & Picture choices, visual expectations, step cards, and start-and-stop signals made routines more visible. & These supports reduced uncertainty and helped children who process spoken directions slowly or need repetition. \\
Sound and movement & Devices, peer conversations, building materials, and moving bodies could create cumulative noise and motion during active lessons. & A child might cover their ears, move away, repeat a question, or stop working because the combined sensory load made it difficult to identify the most important information. \\
Touch and materials & Lessons involved blocks, wires, craft materials, robots, and other objects with varied textures and fine-motor demands. & Material choice could support participation, while an unwanted texture or crowded workspace could become an access barrier unrelated to understanding the academic task. \\
Communication and participation & Students could point, speak, use technology, build, work with a peer, or use an alternative input method. & Flexible participation did not require every child to communicate, move, or respond in the same way. \\
Transitions & Lessons often shifted between whole-group directions, devices, hands-on materials, collaboration, cleanup, and departure. & Rapid or unclear transitions could accumulate into sensory and executive-function overload even when each individual activity was manageable. \\
Regulation and recovery & Movement opportunities and responsive adults were available, but support sometimes depended on an adult recognizing distress in the moment. & Access was less dependable when a child had to become visibly overwhelmed before receiving a break, quieter space, or alternative. \\
\bottomrule
\end{tabular}
\rowcolors{2}{}{}

\reviewsection{Overall Assessment}

I would describe the IGNITE classroom as \textit{meaningfully supportive and partially sensory-friendly}. Its visual supports, flexible participation, movement options, and responsive adults created important access. The strongest support was that participation did not always require one correct body position, communication mode, material, or pace. The laptop example demonstrated that a regulation strategy and an academic task could coexist.

The strongest remaining barrier was cumulative sensory load during active STEM lessons. One Chromebook, robot, conversation, or set of blocks might not be overwhelming by itself. Several devices, peer conversations, moving bodies, visual displays, and tactile materials operating simultaneously could make it difficult for a child to filter information. A second barrier was that available supports were not always visible or self-initiated. The Michigan checklist asks whether visitors are told about busy or noisy periods, whether sensory information is mapped in advance, and whether a quiet space is clearly available \citep{mddi_checklist}. In the classroom, many supports existed, but a child might still need an adult to notice and interpret distress before those supports became available.

This distinction matters because visible behavior can be misread. Covering the ears, keeping a familiar device nearby, leaving a group, avoiding a material, repeating a question, or resisting a transition may communicate overload, uncertainty, or a need for more processing time. These actions do not automatically indicate disinterest, inability, or noncompliance.

\reviewsection{Recommendations}

\textbf{1. Create a visual sensory story and classroom map.} A short, child-friendly guide could show what students may see, hear, touch, and do from entry through dismissal. It could identify the Chromebook area, robotics or building area, supply locations, louder activity zones, seating choices, and the place to take a break. Simple sensory icons could be paired with brief language such as ``devices may make sounds,'' ``materials may feel different,'' and ``you may ask for headphones or a break.'' This would apply the Sensory Detective approach across the child's full experience and respond directly to the checklist's recommendation for sensory maps and advance information.

\textbf{2. Establish a clearly identified low-sensory reset option.} The classroom should have a consistent area where a student can reduce sensory input without being excluded from learning. It could include softer lighting when possible, uncluttered seating, noise-reducing headphones, a visual timer, and a small selection of regulation tools. Students should be taught that using the area is a normal self-advocacy choice rather than evidence of misbehavior. A portable version could be used when a permanent area is unavailable.

\textbf{3. Build sensory choices into each lesson before difficulty appears.} Lesson previews should identify likely sensory demands and offer choices in advance. Students could choose muted device audio, headphones, a quieter seat, a reduced-material workspace, an alternate texture, additional transition time, or a visual version of verbal directions. Teachers could also stagger material distribution and cleanup to reduce crowding and competing noise. Providing these options universally would reduce the need for children to display distress before receiving support.

% \reviewsection{Connection to the Course Theme}

This experience connects to communication as access to agency. Children are often understood through the behavior adults can see, while the environmental conditions producing that behavior remain unexamined. If a student leaves a group, keeps a familiar device nearby, avoids a material, or needs repeated directions, the first question should not be, \textit{``How do we make the student comply?''} A stronger question is, \textit{``What is the student communicating, what sensory or access barrier is present, and what support would expand participation?''}

This also connects to my Puzzle Plan and EQUITAS framework. Shallow evidence records only visible behavior and risks turning an access barrier into a judgment about competence. Richer evidence includes sensory load, communication mode, predictability, student preference, cultural and family knowledge, and environmental design. A sensory-friendly classroom does not ask children to become less complex in order to belong. It redesigns the conditions so that more children can participate, communicate, regulate, and learn with dignity.

\reviewsection{Conclusion}

The sensory walk showed that accessibility is not produced by one tool, one quiet corner, or one responsive adult. It is produced by the relationship among predictability, sensory intensity, communication access, movement choice, and recovery support. The IGNITE classroom already contained meaningful inclusive practices, but those practices would become stronger if sensory information and options were made visible before a child reached overload. The central lesson is that behavior is evidence about the interaction between a person and an environment. Educators improve interpretation when they examine both sides of that interaction and redesign the environment rather than locating the entire problem within the child.

\newpage
\small
\begin{thebibliography}{9}

\bibitem[Solomon R. Guggenheim Foundation(n.d.)]{guggenheim_sensory}
Solomon R. Guggenheim Foundation. (n.d.). \textit{Creating a sensory-friendly environment: Becoming a sensory detective}. Guggenheim for All Toolkit.

\bibitem[Michigan Developmental Disabilities Institute(2025)]{mddi_checklist}
Michigan Developmental Disabilities Institute. (2025). \textit{Sensory-friendly environment checklist}. Wayne State University.

\end{thebibliography}

\end{document}
